¿A qué se dedica un ingeniero forense?

Un ingeniero forense se encarga de investigar el cómo y por qué ocurrió un fallo, accidente o delito. Se especializa en investigar ciberataques, fraudes electrónicos, robos de propiedad intelectual y fugas de datos.

Los ingenieros forenses se encargan de intervenir cuando una organización sufre un ataque para identificar el origen y alcance del mismo, contener la amenaza y evitar la destrucción o alteración de pruebas. Además, utilizan metodologías, hardware y software forense especializado para recolectar datos de un sistema vulnerado y obtener evidencia admisible ante un tribunal; de la misma manera, se encargan de traducir los aspectos técnicos del incidente en informes formales detallados, funcionando como peritos informáticos en procesos judiciales, civiles o penales al exponer la evidencia digital ante jueces, jurados o abogados.
